% Reproduced inherited proof; finding P-F35.
\subsection{Exact indices for the four-variable cubic calibration}

\begin{proposition}[P-F35: both quaternary-cubic indices are twenty-three]
\label{prop:quaternaryindices}
For the ordinary cubic-invariant stretching function of P-F31 with
$n=4$,
\[
 f_4(s)=\dim\left(\operatorname{Sym}^{4s}
                  (\operatorname{Sym}^3\mathbb C^4)\right)^{SL_4},
\]
the strict positivity and saturation indices of GCT VI are exactly
\[
                          p(f_4)=s(f_4)=23.
 \tag{P-F35a}
\]
Its minimal quasipolynomial period is $120$. In particular the
primitive eighth-root poles of P-F31 occur in the classical
explicit series and coexist with these finite indices. This
single dimension gives no asymptotic index bound.
\end{proposition}

\begin{proof}
Salmon's invariant-ring theorem gives algebraically independent
invariants of ordinary degrees $8,16,24,32,40$, together with
one invariant of degree $100$ whose square lies in the polynomial
ring they generate. An exact inspected locator is Huybrechts,
\emph{The geometry of cubic hypersurfaces}, Chapter~4,
Theorem~4.1, printed p.~196; the following proof of Corollary~4.2
also states the resulting graded decomposition. The ordinary
degrees in the polynomial subring are divisible by eight, whereas
those in its degree-$100$ translate are congruent to four modulo
eight. Thus the sum of these two modules is direct; it spans
by the monic quadratic relation. Dividing ordinary degrees by
four gives
\[
 H_4(t)=\sum_{s\ge0}f_4(s)t^s
       =\frac{1+t^{25}}{(1-t^2)(1-t^4)(1-t^6)(1-t^8)(1-t^{10})}.
 \tag{P-F35b}
\]
This also directly checks the nonzero primitive eighth-root
pole: the factor $1-t^8$ vanishes simply there, the other four
factors do not vanish, and $1+\zeta^{25}=1+\zeta\ne0$.

Let $P(y)$ be the global quasipolynomial with generating series
\[
                  \sum_{j\ge0}P(j)z^j
                       =\frac1{\prod_{a=1}^5(1-z^a)}.
\]
We will use its constituent polynomials explicitly. Put
\[
 Q(y)=\frac{y^4}{2880}+\frac{y^3}{96}
            +\frac{31y^2}{288}+\frac{85y}{192}
            +\frac{50651}{86400}.
\]
For an integer residue $r$, define the periodic numbers
\[
 a_r=\begin{cases}1&r\equiv0,1\pmod4,\\-1&r\equiv2,3\pmod4,\end{cases}
 \quad
 b_r=\begin{cases}2&3\mid r,\\-1&3\nmid r,\end{cases}
 \quad
 c_r=\begin{cases}4&5\mid r,\\-1&5\nmid r.\end{cases}
\]
Then the period-$60$ constituent on $r$ is
\[
 P_r(y)=Q(y)+\frac{(-1)^r(2y+15)}{128}
                   +\frac{a_r}{16}+\frac{b_r}{27}+\frac{c_r}{25}.
 \tag{P-F35c}
\]
For completeness, this formula follows by coefficient extraction
from the following exact partial-fraction identity, which can also
be checked by multiplying through by $\prod_{a=1}^5(1-z^a)$:
\begin{align*}
 \frac1{\prod_{a=1}^5(1-z^a)}
 &=\frac1{120(1-z)^5}+\frac1{24(1-z)^4}
   +\frac{31}{288(1-z)^3}+\frac3{16(1-z)^2}
   +\frac{20831}{86400(1-z)}\\
 &\quad+\frac1{64(1+z)^2}+\frac{13}{128(1+z)}
   +\frac{z+1}{16(z^2+1)}
   +\frac{z+2}{27(z^2+z+1)}\\
 &\quad+\frac{z^3+2z^2+3z+4}
                  {25(z^4+z^3+z^2+z+1)}.
 \tag{P-F35d}
\end{align*}
The first line gives $Q$ after using the binomial-series formula.
The two $1+z$ terms give $(-1)^r(2y+15)/128$.
The last three periodic terms give respectively $a_r/16$,
$b_r/27$ and $c_r/25$: multiplication by $1-z$ reduces
the last two numerators to $2-z-z^2$ and
$4-z-z^2-z^3-z^4$. This proves (P-F35c) for every
nonnegative integer argument and hence identifies the whole
constituent polynomials.

Every polynomial $P_r(y-1)$ has nonnegative coefficients.
Indeed its coefficients of degrees four, three and two are
respectively
\[
              \frac1{2880},\qquad\frac{13}{1440},
                       \qquad\frac{113}{1440}.
\]
Its linear coefficient is
\[
             \frac{247}{960}+\frac{(-1)^r}{64}
                     \ \ge\ \frac{29}{120}>0.
\]
Its constant term is
\[
 \frac{20831}{86400}+\frac{13(-1)^r}{128}
                        +\frac{a_r}{16}+\frac{b_r}{27}
                        +\frac{c_r}{25}.
\]
The sum $13(-1)^r/128+a_r/16$ takes the four values
$21/128,-5/128,5/128,-21/128$ on residues modulo four.
Since $b_r\ge-1$ and $c_r\ge-1$, the constant term is at least
\[
 \frac{20831}{86400}-\frac{21}{128}-\frac1{27}-\frac1{25}=0.
 \tag{P-F35e}
\]
All coefficient inequalities have therefore been verified
for every constituent without an interpolation assumption.

Equation (P-F35b) says that the even constituent polynomials
of $f_4$ are of the form $P_r(s/2)$ and the odd ones are
of the form $P_r((s-25)/2)$. To justify the latter globally,
for each odd residue modulo $120$ the equality holds for
all sufficiently large nonnegative $s$; thus it is the
polynomial identity on that residue. The series is a global
quasipolynomial, either by the canonical-ring theorem already
used or directly because its numerator over $(1-t^{120})^5$
has degree $595<600$.

After shifting by $23$, the odd polynomials become
$P_r(s/2-1)$ and therefore have nonnegative coefficients
by (P-F35e) and its preceding calculations. The even
polynomials become $P_r((s+23)/2)$, obtained from the
coefficientwise nonnegative polynomial $P_r(y-1)$ by
substituting $y=(s+25)/2$. A polynomial with nonnegative
coefficients remains such after substitution of an affine
polynomial with nonnegative coefficients. Thus $p(f_4)\le23$.

Every constituent has positive degree-four coefficient
$1/46080$, so none is the zero polynomial. On the other
hand (P-F35b) gives the actual value $f_4(23)=0$:
the denominator contains only even degrees, and the odd
numerator translate first starts at $25$. For any shift
$u<23$, the corresponding nonzero constituent still
vanishes at the positive integer $23-u$. Hence strict
saturation fails for all such shifts, giving
$s(f_4)\ge23$. Coefficientwise positivity implies strict
saturation for every nonzero constituent on positive
integers, so
\[
                   23\le s(f_4)\le p(f_4)\le23.
\]
This proves (P-F35a).

Finally $120$ is a period by the denominator in (P-F35b).
Its actual primitive eighth-root, third-root and fifth-root
poles force every quasipolynomial period to be divisible
by $8,3,5$, respectively. At a primitive third or fifth root,
only the factors of the corresponding divisible degrees
can vanish and the numerator is nonzero ($1+t^{25}$ equals
$1+t$ at a third root and equals two at a fifth root).
Thus their least common multiple $120$ is the minimal period.
\end{proof}

The exploratory exact rational computation is preserved in
\nolinkurl{quaternary_index_calibration.py} and its output file.
The proof above supplies explicit constituent formulas and
uniform inequalities, independently of the finite search.
The exact source URL is
\url{https://www.math.uni-bonn.de/people/huybrech/Notes.pdf}.
